\subsection{Morris 与 Sobol 全局灵敏度模型}

本阶段将 AMR 数量、任务量和瓶颈簇容量作为同一组\textbf{全局灵敏度分析}因素：$Y=f(\theta)=\Cmax(\theta),\theta=(m,n,q_C),$
其中 $m\in\{3,4,5,6\}$，$n\in\{12,13,14,15\}$，$q_C\in\{1,2,3\}$。瓶颈簇 $q_C$ 自动取瓶颈空间簇与容量影子价值分析中排名第一的 C03。连续设计点 $x_i\in[0,1]$ 仅用于抽样，实际求解前按
$
\mathcal{M}_i(x_i)=\ell_{i,r_i},
r_i=1+\min\{L_i-1,\max\{0,\mathrm{round}((L_i-1)x_i)\}\}
$
映射到离散水平 $\mathcal{L}_i=(\ell_{i,1},\ldots,\ell_{i,L_i})$，因此响应为 $Y=f(\mathcal{M}(x))$。

\subsubsection{Morris 基本效应定义}

Morris 方法用\textbf{有限差分基本效应}筛选主导因素。设归一化空间划分为 $p$ 个水平，本实验 $p=4$，步长、第 $r$ 条轨迹上因素 $i$ 的基本效应、对 $R$ 条轨迹汇总定义如下
\[
\Delta=\frac{p}{2(p-1)}=\frac{2}{3}.
\]

\[
EE_i^{(r)}
=\frac{f(\mathcal{M}(x^{(r)}+\Delta e_i))-f(\mathcal{M}(x^{(r)}))}{\Delta}.
\]

\[
\mu_i^\star=\frac{1}{R}\sum_{r=1}^{R}|EE_i^{(r)}|,\qquad
\sigma_i=\sqrt{\frac{1}{R-1}\sum_{r=1}^{R}\left(EE_i^{(r)}-\bar{EE}_i\right)^2}.
\]
其中 $\mu_i^\star$ 表示全局影响强度，$\sigma_i$ 表示影响随背景工况变化的幅度；$\sigma_i$ 较大通常意味着\textbf{非线性效应}或\textbf{交互效应}。

\subsubsection{Sobol 方差分解定义}

Sobol 方法从\textbf{方差分解}度量因素贡献。令 $X=(X_1,\ldots,X_D)$ 且 $D=3$，若 $Y=f(\mathcal{M}(X))$ 平方可积，则
\[
f(X)=f_0+\sum_i f_i(X_i)+\sum_{i<j}f_{ij}(X_i,X_j)+\cdots,\qquad
V=\mathrm{Var}(Y)=\sum_iV_i+\sum_{i<j}V_{ij}+\cdots.
\]
一阶效应和总效应定义为
\[
S_i=\frac{\mathrm{Var}_{X_i}\left(\mathbb{E}_{X_{\sim i}}[Y\mid X_i]\right)}{\mathrm{Var}(Y)},\qquad
S_{T_i}=1-\frac{\mathrm{Var}_{X_{\sim i}}\left(\mathbb{E}_{X_i}[Y\mid X_{\sim i}]\right)}{\mathrm{Var}(Y)}.
\]
$S_i$ 表示因素 $i$ 的独立方差贡献，$S_{T_i}$ 表示包含全部高阶交互后的总贡献。本实验采用 Saltelli/Sobol 设计，$N=64,D=3$，样本规模为 $N(2D+2)=512$。构造 $A,B,A_B^{(i)}$ 后，典型估计量为
\[
\widehat{S_i}=
\frac{\frac{1}{N}\sum_{k=1}^Ny_B^{(k)}(y_{A_B^{(i)}}^{(k)}-y_A^{(k)})}{\widehat V},\qquad
\widehat{S_{T_i}}=
\frac{\frac{1}{2N}\sum_{k=1}^N(y_A^{(k)}-y_{A_B^{(i)}}^{(k)})^2}{\widehat V}.
\]
